\chapter{Aortic Dissection}

\label{ch:ch74_aortic_dissection}

Aortic dissection is a tear in the inner layer of the aorta that allows blood to dissect between the layers of the aortic wall, creating a false lumen. It is a surgical emergency with a mortality rate that increases by 1--2\% per hour if untreated. Rapid, accurate imaging is essential.


\section{1. Patient Perspective: What the Aorta Scan Shows}

If doctors suspect a problem with your aorta, you will have a \textbf{CT angiogram} (CTA) of the chest, abdomen, and pelvis. This scan uses contrast dye and X-rays to produce detailed images of the aorta and its branches.

\begin{itemize}
  \item You will lie on the scanner table; an IV line is placed in the arm.
  \item Contrast dye is injected and the scan is timed precisely to capture the aorta when the dye is at peak concentration.
  \item The scan takes 10--15 seconds. You will be asked to hold your breath.
  \item The scan assesses whether the aortic wall has a tear (dissection flap), whether there is bleeding, and whether the branches supplying your brain, arms, organs, and legs are being affected.
  \item If you cannot receive contrast dye, a non-contrast scan can detect some aortic abnormalities, and MRI or transoesophageal echocardiography may be alternatives.
\end{itemize}


\section{2. Recognition: When to Suspect Dissection}

Aortic dissection should be suspected in any patient with:
\begin{itemize}
  \item Sudden, severe ``tearing'' or ``ripping'' chest or back pain
  \item Blood pressure difference $>$ 20 mmHg between the two arms
  \item New aortic regurgitation murmur
  \item Neurologic deficits (stroke, paraplegia, arm ischemia)
  \item Mesenteric or limb ischemia
  \item Unexplained syncope, shock, or haemopericardium
\end{itemize}

Risk factors include hypertension, Marfan syndrome, Ehlers-Danlos type IV, bicuspid aortic valve, prior aortic surgery, and pregnancy.


\section{3. CT Angiography Protocol}

\begin{itemize}
  \item \textbf{Non-contrast phase}: Detects intramural haematoma (hyperdense crescent in the aortic wall) and helps differentiate from motion artifact.
  \item \textbf{Arterial phase}: Bolus tracking at the ascending aorta (trigger $\sim$150 HU). Visualises the intimal flap, true and false lumens, and branch vessel involvement.
  \item \textbf{Delayed phase} (optional): Assesses end-organ perfusion, delayed contrast extravasation, and stent-graft integrity in known or suspected dissection.
\end{itemize}

\begin{pearlbox}[Timing Matters]
In aortic dissection, the arterial phase must capture the ascending aorta. If the bolus is triggered too late (at the descending aorta), the ascending aorta may be suboptimally opacified, potentially missing a Type A dissection.
\end{pearlbox}


\section{4. Classification Systems}

\subsection{Stanford Classification}
\begin{table}[h]
\centering
\begin{tabular}{@{}p{0.15\textwidth}p{0.42\textwidth}p{0.29\textwidth}@{}}
\hline
\textbf{Type} & \textbf{Involvement} & \textbf{Management} \\
\hline
A & Ascending aorta (with or without descending) & Surgical emergency \\
B & Descending aorta only (distal to left subclavian) & Medical (unless complicated) \\
\hline
\end{tabular}
\caption{Stanford Classification of Aortic Dissection}
\end{table}

\subsection{DeBakey Classification}
\begin{table}[h]
\centering
\begin{tabular}{@{}p{0.20\textwidth}p{0.68\textwidth}@{}}
\hline
\textbf{Type} & \textbf{Description} \\
\hline
I & Originates in ascending aorta, extends to at least the aortic arch \\
II & Confined to ascending aorta \\
IIIa & Extends to descending aorta below diaphragm \\
IIIb & Extends to descending aorta above diaphragm \\
\hline
\end{tabular}
\caption{DeBakey Classification of Aortic Dissection}
\end{table}

\section{5. Key Imaging Findings}

\begin{itemize}
  \item \textbf{Intimal flap}: A linear filling defect separating the true lumen from the false lumen. The true lumen is typically smaller and demonstrates earlier contrast enhancement.
  \item \textbf{True vs false lumen differentiation}: The false lumen is usually larger, has slower contrast opacification, and may show a ``cobweb sign'' (thin strands of intima within the false lumen).
  \item \textbf{Malperfusion}: Branch vessels may be compressed or displaced by the false lumen, causing ischemia to the brain, spinal cord, kidneys, mesentery, or limbs.
  \item \textbf{Complications}: Pericardial effusion/haemopericardium, aortic rupture (mediastinal or pleural haematoma), aortic regurgitation, and coronary artery involvement.
\end{itemize}

\section{6. Acute Aortic Syndromes Spectrum}

Dissection is part of a spectrum of acute aortic syndromes:
\begin{itemize}
  \item \textbf{Intramural haematoma} (IMH): Hemorrhage within the aortic wall without an intimal flap. Appears as a hyperdense crescent on non-contrast CT. May progress to dissection.
  \item \textbf{Penetrating atherosclerotic ulcer} (PAU): Ulceration of an atherosclerotic plaque that penetrates into the media. Often in elderly hypertensive patients with extensive atherosclerosis.
  \item \textbf{Aortic rupture}: Extravasation of contrast or blood outside the aortic wall. Usually catastrophic.
\end{itemize}


\section{7. Medical vs Surgical Management by Type}

\begin{itemize}
  \item \textbf{Stanford Type A} (ascending): Requires \textbf{urgent surgical repair} (open or hybrid). Delay increases mortality.
  \item \textbf{Uncomplicated Stanford Type B}: Managed \textbf{medically} with strict blood pressure and heart rate control (target SBP $<$ 120 mmHg, HR $<$ 60 bpm). Serial imaging surveillance is essential.
  \item \textbf{Complicated Type B} (malperfusion, rapid expansion, rupture, refractory pain): Treated with \textbf{thoracic endovascular aortic repair} (TEVAR) or open surgery.
\end{itemize}

\begin{pearlbox}[Clinical Pearl]
A wide mediastinum on a chest X-ray in a patient with chest pain is a red flag. Although CTA is the gold standard, the initial CXR finding should prompt immediate clinical evaluation and CTA, not delayed investigation.
\end{pearlbox}


\section{8. Structured Reporting Example}

\begin{testbox}[Indications:]
68-year-old male with acute-onset tearing chest pain radiating to the back. Blood pressure 180/110 mmHg in the right arm, 140/90 mmHg in the left. Emergency CTA of the aorta requested.
\end{testbox}
\begin{refbox}[Findings:]
\textbf{Ascending Aorta}: An intimal flap is identified in the ascending aorta, extending from the sinotubular junction to the proximal aortic arch. The true lumen is compressed but patent. No pericardial effusion. No aortic regurgitation.\\
\textbf{Aortic Arch \& Descending Aorta}: The dissection flap extends into the aortic arch and descending thoracic aorta, terminating at the level of the coeliac trunk. The false lumen is patent and larger than the true lumen. No contrast extravasation.\\
\textbf{Branch Vessels}: The left common carotid and left subclavian arteries arise from the true lumen. The coeliac trunk and superior mesenteric artery are perfused by the true lumen. No malperfusion.\\
\textbf{Abdominal Aorta}: No extension below the renal arteries. Both renal arteries arise from the true lumen with normal enhancement of both kidneys.
\end{refbox}
\begin{pearlbox}[Impression:]
1. Stanford Type A (DeBakey Type I) aortic dissection involving the ascending aorta, arch, and proximal descending thoracic aorta.
2. No pericardial effusion, aortic regurgitation, or branch vessel malperfusion.
3. Urgent cardiothoracic surgical consultation recommended.
\end{pearlbox}

\clearpage
